% Project abstract (max. 4 pages), Web Data Integration HWS 2026
% Springer LNCS template: llncs.cls v2.26. Build: latexmk -pdf proposal.tex
\documentclass[runningheads,a4paper]{llncs}
\usepackage[T1]{fontenc}
\usepackage{booktabs}
\usepackage{tabularx}
\usepackage{float} % [H] keeps each table between two paragraphs instead of inside a sentence
\setlength{\intextsep}{9pt plus 2pt minus 2pt}
\usepackage[hidelinks]{hyperref}
\urlstyle{same}
% let long links break after hyphens as well
\makeatletter\g@addto@macro{\UrlBreaks}{\do\-}\makeatother
\newcommand{\mv}{\,(MV)}
\newcolumntype{L}{>{\raggedright\arraybackslash}X}
\newcolumntype{P}[1]{>{\raggedright\arraybackslash}p{#1}}

\begin{document}

\title{Integrating Worldwide University Data from Wikidata, DBpedia, ROR and THE}
\subtitle{Project Abstract -- Web Data Integration, HWS 2026}
\titlerunning{Integrating Worldwide University Data}

\author{Laith Sandouk \and Mohammad Hadad \and Sanjeevani Rajpurohit \and Chellapuram Keerthan Reddy
\and Lakpa Dolma Sherpa}
\authorrunning{Web Data Integration, HWS 2026}
\institute{University of Mannheim, Germany}

\maketitle

\section{Use Case}

Students, researchers and other stakeholders who need information about universities worldwide
face a fragmented landscape: registries, knowledge bases and rankings differ in format,
definitions and coverage, which makes institutions hard to compare. The Research Organization
Registry (ROR) has reliable names and locations, but few descriptive facts. Wikidata and
DBpedia add founding years, student numbers and institution types, but they are incomplete and
often contradict each other. The Times Higher Education (THE) World University Rankings add
scores, student statistics and subjects, but for only 3,100 institutions.

We integrate these four sources into one coherent dataset that describes every institution
exactly once. Students can compare universities by location, ranking and subjects; analysts,
policy makers and university administrators can benchmark institutions without consolidating
sources by hand.

\section{Datasets: Schema and Basic Profile}

Table~\ref{tab:datasets} profiles the four datasets, one more than required. Each contains a
single class, \emph{University}, with one record per entity.

\paragraph{Data sources and preprocessing.}
Our Python scripts retrieved all data on 2 October 2026 and store each source as one CSV file
with one row per university. Attributes with several values are kept as lists, and website
links that differ only in \texttt{http}/\texttt{https} or \texttt{www} are merged.
\begin{itemize}
\setlength{\emergencystretch}{2em}% lets the long links break without sticking out
\item \emph{Wikidata}, \url{https://query.wikidata.org/sparql}: all items that are an instance
of \emph{university} (Q3918) or of one of its subclasses. Identifiers of countries, places and
types were replaced by English labels; per university we kept the earliest founding year, the
preferred or most recent student number and the city of its headquarters or location.
\item \emph{DBpedia}, \url{https://dbpedia.org/sparql}: all resources of class
\texttt{dbo:University}, fetched with one paged query per attribute, because the endpoint
returns at most 10,000 rows. We kept the earliest founding year and the largest student number,
cut leftover wiki markup off the website links and unwrapped archived links.
\item \emph{ROR}, data dump v2.13 of 22 September 2026,
\url{https://doi.org/10.5281/zenodo.22902037}: of its 141,528 records we kept the 30,131 of
type \emph{education} that are not withdrawn (withdrawn records are merged duplicates).
\item \emph{THE},
\url{https://www.timeshighereducation.com/world-university-rankings/2026/world-ranking}: the
table that the web page loads, with 2,191 ranked institutions and 927 ``reporters'' that
supplied data but got no rank. We added the ranking status, turned text such as ``22,005'' and
``43\%'' into numbers and split the subjects into a list.
\end{itemize}

\begin{table}[H]
\caption{Datasets}\label{tab:datasets}
\footnotesize
\setlength{\tabcolsep}{3pt}
\begin{tabularx}{\textwidth}{@{}lP{1.5cm}lrrL@{}}
\toprule
Dataset & Source (*) & Format &
\multicolumn{1}{P{1.1cm}}{\raggedleft\# of entities} &
\multicolumn{1}{P{0.75cm}}{\raggedleft\# of attr.} & List of attributes (**)\\
\midrule
Wikidata & SPARQL endpoint & CSV & 22,697 & 16 &
Wikidata id, name, alt.\ names\mv, acronyms\mv, country, country code, city, website, founding
year, latitude\mv, longitude\mv, students\mv, types, member of\mv, ROR id\mv, Wikipedia
title\mv\\
\addlinespace[2pt]
DBpedia & SPARQL endpoint & CSV & 27,728 & 14 &
DBpedia id, name, alt.\ names\mv, country, state\mv, city, website, founding year\mv,
latitude\mv, longitude\mv, students\mv, types\mv, affiliations\mv, Wikidata id\\
\addlinespace[2pt]
ROR & Data dump (Zenodo) & CSV & 30,131 & 18 &
ROR id, name, alt.\ names, acronyms\mv, country, country code, region, city, website,
domains\mv, founding year, latitude, longitude, types, status, parent organisation\mv,
Wikidata id, Wikipedia title\mv\\
\addlinespace[2pt]
THE & Ranking web page & CSV & 3,118 & 22 &
THE id, name, country, ranking status, rank, overall score, 5 pillar scores, 5 pillar ranks,
students, students per staff, international students \%, female students \%, subjects, profile
URL\\
\bottomrule
\end{tabularx}

\smallskip
(*) See ``Data sources and preprocessing''. \enspace (**) (MV) = more than 30\% missing values.
The identifiers and the links to other sources serve only as gold standard for the matching.
\end{table}

\paragraph{Missing values.}
Over all attributes, 32\% of the values are missing in Wikidata, 33\% in DBpedia, 17\% in ROR
and 17\% in THE. This does not hinder
identity resolution, because the attributes we match on are well filled: the name is never
missing and the country only in 3\% of the Wikidata and 19\% of the DBpedia records. For the
fusion the gaps are the reason for the project, because the sources fill each other's gaps:
student numbers are missing in 84\% of the Wikidata and 67\% of the DBpedia records but in only
1\% of the THE records, and the coordinates missing in 31\% of the Wikidata records are complete
in ROR.

\paragraph{Ranking coverage.}
The largest gap is the ranking: THE lists only 3,118 institutions, 6\% of the joint dataset
(2,191 with a rank and 927 reporters without one). We see two options and would welcome
feedback: (a) keep all institutions and label those that THE does not list as \emph{unranked}
in the ranking status; or (b) restrict the final dataset to the institutions that THE lists.
Option (b) makes the dataset dense and, with 3,118 entities, still meets the required 2,500, but
it stays below the range of more than 10,000 that the requirements call good.

\section{Integrated Schema and Overlap with Input Schemata}

The integrated schema has one class, \emph{University}, with 34 attributes
(Table~\ref{tab:schema}). 14 attributes occur in at least two datasets and nine in at least
three, eight of them besides the name, so fusion by voting is possible. Six attributes are
lists. No single attribute serves as a key: the identifiers of the sources are not part of the
schema, and the name alone is ambiguous (in ROR, ``National University'' exists in four
countries). A university is identified by the combination of name and country, which is unique
for 97.2\% of the Wikidata, 99.8\% of the DBpedia and ROR, and all THE records; the city
separates most of the rest.

\begin{table}[H]
\caption{Attribute Intersection with Integrated Schema}\label{tab:schema}
\footnotesize
\begin{tabularx}{\textwidth}{@{}P{4.9cm}lL@{}}
\toprule
Attribute name & Attribute type & Datasets in which the attribute is found\\
\midrule
name, country & string & Wikidata, DBpedia, ROR, THE\\
alternative names & list of strings & Wikidata, DBpedia, ROR\\
city & string & Wikidata, DBpedia, ROR\\
website & URL & Wikidata, DBpedia, ROR\\
founding year & date (year) & Wikidata, DBpedia, ROR\\
latitude, longitude & decimal & Wikidata, DBpedia, ROR\\
number of students & integer & Wikidata, DBpedia, THE\\
acronyms & list of strings & Wikidata, ROR\\
country code (ISO) & string & Wikidata, ROR\\
region (DBpedia: state) & string & DBpedia, ROR\\
institution types, memberships & list of strings & Wikidata, DBpedia\\
parent organisation, status & string & ROR\\
internet domains & list of strings & ROR\\
ranking status, rank, overall score & string & THE\\
5 pillar scores, 5 pillar ranks & decimal, integer & THE\\
students per staff, international students~\%, female students~\% & decimal & THE\\
subjects offered & list of strings & THE\\
\bottomrule
\end{tabularx}
\end{table}

The mapping is not one-to-one everywhere: the \emph{affiliations} of DBpedia also contain parent
universities, which we separate during translation, and the \emph{types} of ROR only hold
registry categories such as \emph{education} or \emph{funder}, so they are not mapped to the
institution types.

\section{Entities Contained in Multiple Datasets}

\paragraph{Wikidata, DBpedia and ROR.}
These three sources reference each other through ROR identifiers, Wikidata identifiers and
Wikipedia article titles (the DBpedia identifiers); 59\% to 93\% of the records of each source
carry such a link. Following these links through the complete files, instead of a small
sample, groups the records into 50,036 distinct institutions; 19,895 of them are contained in
at least two sources and 9,829 in all three.

\paragraph{THE.}
THE shares no identifier with the other sources. Comparing normalised names with the names and
alternative names of the other sources links 2,626 of the 3,118 THE universities (84\%). This
is a lower bound: of a random sample of 100 that we checked by hand, 98 are contained in another
source, 16 of them under a different spelling. With THE, the joint dataset describes about
50,500 institutions: 19,993 in at least two sources, 10,125 in at least three and 2,225 in all
four (Table~\ref{tab:requirements}).

\paragraph{Identity resolution remains a real task.}
The identifier columns are removed from the matching input and serve only as gold standard;
for THE it is labelled by hand. Names alone do not suffice: 25\% of the linked Wikidata--ROR
and 18\% of the linked Wikidata--DBpedia pairs differ in name even after normalisation.

\begin{table}[H]
\caption{Project requirements and how the project meets them}\label{tab:requirements}
\footnotesize
\setlength{\tabcolsep}{4pt}
\begin{tabularx}{\textwidth}{@{}llL@{}}
\toprule
& Requirement & Our project\\
\midrule
1. & 3 different datasets & 4: Wikidata, DBpedia, ROR, THE\\
2. & $\geq$\,2,500 entities in total, good: $>$\,10,000 & about 50,500 (from 83,674 records)\\
3. & $\geq$\,1,000 entities in at least two datasets & about 20,000\\
4. & $\geq$\,8 attributes, identifiable by two of them & 34; name + country, no identifier\\
5. & $\geq$\,5 attributes in two datasets, some in three & 14 in two, 9 in three or more\\
6. & At least one list attribute & 6, e.g.\ alternative names, subjects\\
\bottomrule
\end{tabularx}
\end{table}

\paragraph{Independence of the sources.}
ROR is curated by its own team, and the THE figures are submitted by the universities
themselves. DBpedia is extracted from the English Wikipedia, on which Wikidata partly draws as
well, so we count the two as one source for name and coordinates (63\% of the linked pairs have
identical coordinates). For the attributes fused by voting they behave independently: where
exactly one of the three linked sources deviates in the founding year, Wikidata and DBpedia
agree against ROR least often (5.6\%, against 8.5\% and 6.6\% for the other pairings). Without
DBpedia, only the country would remain in three datasets.

\end{document}
